\subsection*{NT-E12. Distinct CM fields obstruct simultaneous integral extremes}

Let \(E_1,E_2\) be elliptic curves over a number field \(K\),
with all their complex multiplication defined over \(K\).
Suppose their imaginary quadratic CM fields \(M_1,M_2\)
are distinct. At a prime of common good reduction, put
\[
 q=N_K\mathfrak p,\qquad
 a_i=q+1-\#E_i(\mathcal O_K/\mathfrak p).
 \]
Then the following implication is unconditional:
\[
            a_1^2=a_2^2
       \quad\Longrightarrow\quad
       q\text{ is a square and }a_1^2=a_2^2=4q.
                                                               \tag{E12.1}
\]
In particular, at an absolute-degree-one prime the two
absolute traces are different. Simultaneous integral
extremal events
\[
            |a_1|=|a_2|=\lfloor2\sqrt q\rfloor
                                                               \tag{E12.2}
\]
can occur only at square norms and exact Hasse endpoints.
For all primes of norm at most \(x\) satisfying (E12.2),
the elementary bound
\[
        \#\{\mathfrak p:\text{(E12.2)},\ N_K\mathfrak p\le x\}
                              \le [K:\mathbf Q]\sqrt x
                                                               \tag{E12.3}
\]
holds. No RH assumption is used.

\begin{proof}
We give the centralizer argument, including the rational
Frobenius exception. Let \(E/K\) be either curve, with
CM field \(M\), and let \(\overline E\) be its good
reduction over \(\mathbf F_q\). Good-reduction
specialization injects \(\operatorname{End}_K(E)\) into
\(\operatorname{End}_{\mathbf F_q}(\overline E)\).
Thus \(M\) embeds in
\(A=\operatorname{End}_{\overline{\mathbf F}_q}^0(\overline E)\).
Every element in that embedded field commutes with the
\(q\)-power Frobenius \(\pi\), since its endomorphisms
are defined over \(\mathbf F_q\).

Choose any auxiliary prime \(\ell\ne\operatorname{char}\mathbf F_q\).
The faithful Tate-module action injects
\(A\otimes\mathbf Q_\ell\) into
\(\operatorname{End}_{\mathbf Q_\ell}(V_\ell\overline E)\),
whose underlying space has dimension two.
The action of the two-dimensional \'{e}tale algebra
\(M\otimes\mathbf Q_\ell\) is faithful and has rank one
over that algebra. If the algebra is a field, this
follows from the two dimensions; if it is
\(\mathbf Q_\ell\times\mathbf Q_\ell\), faithfulness
makes each of its two idempotent summands nonzero,
so each has dimension one.
Its centralizer in the two-by-two endomorphism algebra
is therefore exactly \(M\otimes\mathbf Q_\ell\).

Commutation with a basis of \(M\) is a finite system of
rational linear equations in \(A\). The centralizer of
\(M\) in \(A\) consequently has rational dimension at
most two, and already contains \(M\). It equals \(M\).
Since \(\pi\) commutes with \(M\), we conclude \(\pi\in M\).
These specialization and faithful-action facts are
the same classical good-reduction interface used in
the full proof of NT-E10; no assertion about a
Tate-conjecture surjectivity is needed.

The characteristic polynomial on \(V_\ell\) is
\[
                       X^2-aX+q .
 \]
If \(\pi\notin\mathbf Q\), then \(\mathbf Q(\pi)=M\);
its minimal polynomial has degree two and must be the
displayed characteristic polynomial. In particular
\[
                 M=\mathbf Q(\sqrt{a^2-4q}),
                         \qquad a^2-4q<0.               \tag{E12.4}
 \]
If instead \(\pi=t\in\mathbf Q\), its action is scalar.
Its trace and determinant give \(a=2t\) and \(q=t^2\).
As \(q\) is an integer, a rational square root \(t\)
is an integer. Thus \(q\) is a square and \(a^2=4q\).
Conversely a repeated root of the displayed polynomial
forces \(\pi=a/2\): inside the field \(M\), an element
whose square \((\pi-a/2)^2\) is zero is zero.

Now suppose \(a_1^2=a_2^2\). If this common value is
not \(4q\), both Frobenius elements are nonrational,
and (E12.4) identifies both \(M_1\) and \(M_2\) with
the same imaginary quadratic field. This contradicts
their distinctness. Hence the common value is \(4q\),
which proves (E12.1). If \(q=p\) is a rational prime
it is not a square, giving the degree-one conclusion.
More generally every odd absolute residue degree is
excluded.

Event (E12.2) implies equality of squared traces, so
it can occur only as stated. At square norm \(q=p^{2f}\le x\),
the underlying rational prime satisfies \(p\le\sqrt x\).
There are at most \([K:\mathbf Q]\) prime ideals of \(K\)
over each rational prime. Summing this bound over
the rational primes at most \(\sqrt x\) proves
(E12.3), even with the weaker count by all integers.
\end{proof}


\paragraph{Earlier rational special case.}
Bucur--Fit\'e--Kedlaya, Remark~4.11, published p.~1743
(author v2 p.~22), already gives the obstruction
$a_p(E_1)=a_p(E_2)=\lfloor2\sqrt p\rfloor$ over
$\mathbf Q$ for distinct CM fields and $p>3$,
crediting Andrew Sutherland. The formulation here
retains arbitrary base fields, equal absolute traces,
the precise square-norm exception, and the quantitative
count of higher-degree primes. The rational positive-extreme
special case is not claimed as a new literature result.

\paragraph{Meaning of the shrinking-boundary gate.}
Each separate extremal CM trace asks for an angle
of width comparable to \(q^{-1/4}\). Two such
requirements have total shrinking exponent \(1/2\),
exactly the boundary excluded from NT-A16.
Equations (E12.1)--(E12.3) exhibit a genuine integral
Frobenius obstruction there: continuous torus Haar
mass cannot be used to infer degree-one prime
infinitude at this scale. The obstruction concerns
the two integral traces and their quadratic
endomorphism algebras; it is compatible with
equidistribution in every fixed positive-volume box.
It persists after imposing any finite Frobenius
condition, since that only removes primes.
Geometrically the scalar commutant of each CM action
determines the nonrational Frobenius field, so two
distinct CM algebras cannot give the same noncentral
characteristic polynomial.
Full proof: \texttt{cm\_joint\_endpoint\_obstruction.tex},
NT-E12, also \texttt{section.tex}; run
\texttt{GRH\_NIGHT/number\_theory}.
